\documentclass{article}

\usepackage{tabularx}
\usepackage{booktabs}

\title{Problem Statement and Goals\\\progname}

\author{\authname}

\date{}

\input{../Comments.text}
\input{../Common.text}

\begin{document}

\maketitle

\begin{table}[hp]
\caption{Revision History} \label{TblRevisionHistory}
\begin{tabularx}{\textwidth}{llX}
\toprule
\textbf{Date} & \textbf{Developer(s)} & \textbf{Description of Changes}\\
\midrule
Sept 25th 2026 & Ayush Patel & Added custom documents\\
... & ... & ...\\
\bottomrule
\end{tabularx}
\end{table}

\section{Problem Statement}

\wss{You should check your problem statement with the
\href{https://smiths.github.io/capTemplate/Checklists/ProbState-Checklist.pdf}
{problem statement checklist}}. 

\wss{You can change the section headings, as long as you include the required
information.}

\subsection{Problem}

\wss{What problem are you trying to solve?  Why is it important? Do not talk
about how you will solve the problem, only what the problem is.  AI is rarely
part of the problem.  The problem is ``what'' you want to do, not ``how'' you
want to do it.  The problem statement should be written in a way that is
understandable to a non-technical audience.}

\subsection{Inputs and Outputs}

\wss{Characterize the problem in terms of ``high level'' inputs and outputs.  
Use abstraction so that you can avoid details.}

\subsection{Stakeholders}

\subsection{Environment}

\wss{Hardware and Software Environment for the users.  The developer environment
is summarized as part of the developer plan.}

\section{Goals}

\section{Stretch Goals}

\wss{Goals you hope to get to, but not necessary.  They are also helpful if the
scope of the original project needs to be expanded.}

\section{Custom Documents (CDs)}

\wss{For SE Capstone: List your custom documents (CDs).  Potential CDs include
usability testing, code walkthroughs, user documentation, formal proof,
GenderMag personas, Design Thinking, etc.  (The full list is on the course
outline and in Lecture 02.) Normally the number of CDs will be two (one per
term).  Approval of the CDs will be part of the discussion with the instructor
for approving the project. The CDs, with the approval (or request) of the
instructor, can be modified over the course of the term.}

\begin{itemize}
    \item \textbf{User/Stakeholder/Client Interview Report} -- Documents the
    interviews conducted with users, stakeholders, and clients to gather
    requirements for the application.

    \item \textbf{Usability Testing Report} -- Documents the usability testing of the
    application, including the goals of the testing, the specific tasks we ask users to complete, the specific questions we ask users, and the results
\end{itemize}


\newpage{}

\section*{Appendix --- Reflection}

\wss{Not required for CAS 741}

\input{../Reflection.text}

\textbf{Waqar Ul-Hassan:}

\begin{enumerate}
    \item What went well while writing this deliverable? 
    
    The team has been awesome in picking up work and I think that's been the best part of it so far. We haven't had any issues with assigning work and infact most of us have been trying to get more work to make sure that we are all keeping up. Furthermore, being able to get to know my team members are bit better as I didn't know two of them before this project has been nice as well. We also had to drop our original project	we had committed to and I think our team recovered from that pretty, I was a bit worried initially but we got everything together quite quickly!

    \item What pain points did you experience during this deliverable, and how
    did you resolve them?

    I had a couple pain points through this process, the first was of course the fact that our first project ended up being somewhat complete already which was something that caught us off-guard. I Feel as though we did a great job in resolving this by communicating with our supervising professor as well as Dr. Smith and then ensuring we had a backup project incase we wanted to switch. Another Pain point I had was when deciding some of the metrics for how to track all the work being done equally. The solution to this was just holding discussions, we discussed metrics like commits, number of issues, and lines of code. Overall, we ended up deciding on the number assigned issues as the metric because we felt we could best distribute that one. 

    \item How did you and your team adjust the scope of your goals to ensure
    they are suitable for a Capstone project (not overly ambitious but also of
    appropriate complexity for a senior design project)?

    We initially had some more particular problem solutions for our project which we faced personally in finding a supervisor, projects, and navigating through all the information in general. After the discussion with our TA Chris I realized that maybe we needed one or two more smaller things that we should aim to reach as part of our main goal. With this in mind we had discussions with each other and decided on what our goals and stretch goals are now. I also realize as we meet with stakeholders these requirements could change and our TA had said that in the case that they do, we can discuss it with him to update our goals to ensure everything makes sense.

\end{enumerate}  

\end{document}